\begin{helper}
Let $N, n, m$ be natural numbers with $n \le N$ and $m \le N$. Let $Y$ be a hypergeometric random variable counting the number of marked elements in a sample of size $n$ drawn uniformly without replacement from a population of size $N$ containing $m$ marked elements. Let $q = m/N$, with the convention \(q=0\) when \(N=0\), and let $\mu = nq$. For any $\lambda \in \mathbb{R}$, the moment generating function satisfies
\[
\mathbb{E}[e^{\lambda Y}] \le (1 - q + q e^\lambda)^n \le \exp(\mu(e^\lambda - 1)).
\]
\end{helper}
\begin{proof}
If $N=0$, then $n \le 0 \implies n=0$ and $m \le 0 \implies m=0$.  By the stated convention, equivalently by the real-division convention used in the Lean statement, \(q=(m:\mathbb R)/(N:\mathbb R)=0\), and hence \(\mu=nq=0\). The sample sum is exactly $0$ with probability $1$, so $\mathbb{E}[e^{\lambda Y}] = 1$. The right hand sides are $(1-0+0)^0 = 1$ and $\exp(0) = 1$, so the inequalities hold. Assume henceforth $N \ge 1$.

Let the population be $M \subseteq U$ with $|U|=N$ and $|M|=m$. Let $v \in \mathbb{R}^N$ be the indicator vector of the population, where $v_k = 1$ if $k \in M$ and $0$ otherwise.
The without-replacement sample sum $Y$ has the same distribution as $\sum_{i=1}^N b_i v_{\pi(i)}$, where $\pi$ is a uniformly random permutation in the symmetric group $S_N$, and $b \in \mathbb{R}^N$ is the vector $b = (1, \dots, 1, 0, \dots, 0)$ containing $n$ ones and $N-n$ zeros.
The expected value is $\mathbb{E}[e^{\lambda Y}] = \frac{1}{N!} \sum_{\pi \in S_N} \exp(\lambda \sum_{i=1}^N b_i v_{\pi(i)})$.

Let $Z$ be the sum of $n$ independent draws with replacement from the population. $Z$ can be generated by drawing a sequence $X_1, \dots, X_n$ independently and uniformly from $\{1, \dots, N\}$, and computing $\sum_{j=1}^n v_{X_j}$. Let $c_i$ be the number of times $i$ appears in the sequence $X$. Then $Z = \sum_{i=1}^N c_i v_i$.
By symmetry, the distribution of $Z$ is unchanged if we apply a uniformly random permutation $\pi \in S_N$ to the population. Thus,
\[
\mathbb{E}[e^{\lambda Z}] = \mathbb{E}_X \left[ \frac{1}{N!} \sum_{\pi \in S_N} \exp\left(\lambda \sum_{i=1}^N c_i v_{\pi(i)}\right) \right].
\]
For any fixed sequence $X$, the count vector $c = (c_1, \dots, c_N)$ consists of non-negative integers summing to $n$. Therefore, the sum of the largest $k$ elements of $c$ is at least $\min(k, n)$. The vector $b$ has exactly $n$ ones, so the sum of its largest $k$ elements is exactly $\min(k, n)$. This implies that $c$ majorizes $b$.

We now prove, in this finite situation, the convex-hull step usually called
the Hardy-Littlewood-P\'olya theorem.  Sort the count vector in nonincreasing
order and call the sorted vector \(x\).  The vector \(b=(1,\ldots,1,0,\ldots,0)\)
is already nonincreasing.  The majorization inequalities are
\[
S_r(x):=\sum_{h=1}^r(x_h-b_h)\ge0\qquad(1\le r\le N),
\]
and the total-sum equality is \(S_N(x)=0\).

If \(x=b\), then \(b\) is a permutation of the original count vector \(c\).
Otherwise let \(i\) be the first index with \(x_i\ne b_i\).  The prefix
inequality at \(i\) forces \(x_i>b_i\).  Since the total sums agree, there is
some \(j>i\) with \(x_j<b_j\); take the first such \(j\).  For every
\(i\le r<j\), the slack \(S_r(x)\) is strictly positive: before \(j\) no
negative discrepancy has appeared after the first positive discrepancy at
\(i\).  Define
\[
\delta=\min\Bigl(x_i-b_i,\ b_j-x_j,\ \min_{i\le r<j}S_r(x)\Bigr)>0.
\]
All quantities in this minimum are positive integers, so \(\delta\) is a
positive integer and the next vector still has integer entries.
Let \(y\) be obtained from \(x\) by replacing \(x_i\) by \(x_i-\delta\) and
\(x_j\) by \(x_j+\delta\), leaving all other coordinates fixed.  The total
sum is unchanged.  The prefix slacks of \(y\) are the same as those of \(x\)
except for \(i\le r<j\), where
\[
S_r(y)=S_r(x)-\delta\ge0
\]
by the definition of \(\delta\).  Thus \(y\) still majorizes \(b\) in the
original order.  If \(y^\downarrow\) denotes the nonincreasing rearrangement
of \(y\), then \(y^\downarrow\) also majorizes \(b\), because among all
permutations of \(y\), the nonincreasing one maximizes every initial segment
sum.

The transfer vector \(y\) lies in the convex hull of the permutations of
\(x\).  Indeed \(x_i>x_j\): otherwise \(x_i>b_i\ge b_j>x_j\) would fail.
Set \(\theta=\delta/(x_i-x_j)\).  Since
\[
x_i-x_j=(x_i-b_i)+(b_i-b_j)+(b_j-x_j)
\]
and the three terms on the right are nonnegative except that the first and
last are positive, the chosen \(\delta\) satisfies \(0\le\theta\le1\).
On coordinates \(i,j\),
\[
(x_i-\delta,\ x_j+\delta)
=(1-\theta)(x_i,x_j)+\theta(x_j,x_i),
\]
and all other coordinates are unchanged.  Hence \(y\), and therefore its
rearrangement \(y^\downarrow\), belongs to the convex hull of the
permutations of \(x\).

It remains to justify termination.  Use the integer measure
\[
D(x)=\sum_{h=1}^N |x_h-b_h|.
\]
Before resorting, the transfer decreases the two affected absolute-value
terms by \(\delta\) each, since \(x_i>b_i\) and \(x_j<b_j\) and
\(\delta\le x_i-b_i,\ b_j-x_j\).  Thus \(D(y)=D(x)-2\delta\).  Resorting
does not increase the distance to the already sorted vector \(b\): swapping
an inverted pair \(y_r<y_s\) with \(r<s\) cannot increase
\(|y_r-b_r|+|y_s-b_s|\), because \(b_r\ge b_s\) and the absolute-value cost
has the usual Monge inequality
\[
|a-c|+|d-e|\le |a-e|+|d-c|
\quad\text{whenever }a\le d\text{ and }c\le e .
\]
Applying this adjacent-swap comparison until the vector is sorted gives
\(D(y^\downarrow)\le D(y)<D(x)\).  The measure \(D\) is a nonnegative
integer, so the iteration terminates.  Its only terminal majorizing sorted
state is \(b\); otherwise the index choice above would produce another
positive transfer.  Composing the convex-combination inclusions from the
successive steps shows that \(b\) lies in the convex hull of the permutations
of the original count vector \(c\).

Therefore there exist non-negative weights \(w_\sigma\), indexed by
permutations \(\sigma\in S_N\), whose sum is \(1\), such that
\[
b=\sum_{\sigma\in S_N}w_\sigma\,\sigma(c),
\]
where \(\sigma(c)\) is the vector with entries \(c_{\sigma^{-1}(i)}\).
Consider the function $g(x) = \frac{1}{N!} \sum_{\pi \in S_N} \exp(\lambda \sum_{i=1}^N x_i v_{\pi(i)})$. The exponential function is convex, so $g$ is convex. Furthermore, $g$ is symmetric, meaning $g(\sigma(c)) = g(c)$ for any $\sigma \in S_N$.
Applying Jensen's inequality to the convex combination, we obtain
\[
g(b) = g\left(\sum_{\sigma \in S_N} w_\sigma \sigma(c)\right) \le \sum_{\sigma \in S_N} w_\sigma g(\sigma(c)) = \sum_{\sigma \in S_N} w_\sigma g(c) = g(c).
\]
Therefore, for any sequence $X$, we have $g(b) \le g(c)$. Taking the expectation over all sequences $X$ gives $\mathbb{E}[e^{\lambda Y}] = g(b) \le \mathbb{E}_X[g(c)] = \mathbb{E}[e^{\lambda Z}]$.

The sum $Z$ has a binomial distribution with parameters $n$ and $q = m/N$. Its moment generating function is
\[
\mathbb{E}[e^{\lambda Z}] = (1 - q + q e^\lambda)^n.
\]
Using the inequality $1 + x \le e^x$, we have $1 - q + q e^\lambda = 1 + q(e^\lambda - 1) \le \exp(q(e^\lambda - 1))$. Raising this to the power $n$ yields
\[
\mathbb{E}[e^{\lambda Z}] \le \exp(nq(e^\lambda - 1)) = \exp(\mu(e^\lambda - 1)).
\]
Combining these inequalities yields $\mathbb{E}[e^{\lambda Y}] \le (1 - q + q e^\lambda)^n \le \exp(\mu(e^\lambda - 1))$, completing the proof.
\end{proof}
